\chapter{C-Peptide}

\section*{What Is This Test?}
The C-peptide test measures the amount of C-peptide in the blood. C-peptide (connecting peptide) is a 31-amino-acid polypeptide released in equimolar amounts with insulin from pancreatic beta cells. Proinsulin is cleaved in the beta cell granules by prohormone convertases PC1/3 and PC2, producing insulin and C-peptide, plus two dipeptides (Arg31, Arg32). Measurement of C-peptide is the most accurate practical way to estimate endogenous insulin secretion because C-peptide is not extracted by the liver (hepatic extraction of insulin is 50--60\%), has a longer half-life (30--40 minutes vs. 4--6 minutes for insulin), and is not affected by hemolysis (which degrades insulin), insulin antibodies, or exogenous insulin administration. It is essential for classifying diabetes, evaluating fasting hypoglycemia, and monitoring beta-cell function in patients on insulin therapy.

\section*{Alternative Names}
\begin{itemize}
  \item Connecting Peptide
  \item Insulin C-Peptide
  \item C-Peptide Level
  \item C-Peptide Test
\end{itemize}
\section*{Principle}
C-peptide is measured by immunometric methods (radioimmunoassay [RIA], chemiluminescent immunoassay [CLIA], ELISA, or electrochemiluminescence [ECLIA]). The assay uses antibodies that are highly specific for C-peptide and do not cross-react with insulin or proinsulin. Fasting C-peptide levels reflect basal endogenous insulin secretion. The stimulation C-peptide test (C-peptide 6 minutes after 1 mg glucagon IV) assesses maximal beta-cell secretory capacity. This is the most robust measure of beta-cell function. When there is residual endogenous insulin secretion, C-peptide is present in plasma and urine; when beta cells are destroyed (type 1 diabetes), C-peptide levels are very low or undetectable. Urine C-peptide (measured as C-peptide-to-creatinine ratio) correlates well with the mean plasma C-peptide over 24 hours and can be used when repeated blood sampling is impractical (e.g., in children).

\section*{History}
C-peptide was discovered by Donald Steiner and colleagues in 1967 at the University of Chicago during the characterization of proinsulin biosynthesis \cite{steiner1967insulin}. The discovery revealed that insulin is synthesized as a single-chain precursor (proinsulin) that is cleaved into insulin and C-peptide. This established the structural basis of insulin biosynthesis and earned Steiner lasting recognition. Measurement of C-peptide became clinically available in the 1970s. It was used extensively in the Diabetes Control and Complications Trial (DCCT, 1983--1993), which demonstrated that sustained endogenous C-peptide secretion is associated with improved glycemic control and reduced risk of complications. The glucagon-stimulated C-peptide test was established as the gold standard for assessing residual beta-cell function in type 1 diabetes. Contemporary guidelines continue to recommend C-peptide measurement for the classification of diabetes when the type is uncertain, particularly in adults with atypical presentations \cite{cryer2009evaluation}.

\section*{Why This Test Is Ordered}
\begin{itemize}
  \item Classification of diabetes mellitus (type 1 vs. type 2): a C-peptide <0.2 nmol/L (or undetectable) with hyperglycemia strongly suggests type 1 diabetes (T1D) or latent autoimmune diabetes in adults (LADA). A C-peptide >0.3--0.6 nmol/L with hyperglycemia is more consistent with type 2 diabetes (T2D), with the caveat that C-peptide is low in type 1 diabetes but can also be low in long-standing type 2 diabetes with marked beta-cell exhaustion (``type 2 diabetes with insulin deficiency'')
  \item Assessment of residual beta-cell function in established type 1 diabetes: monitoring C-peptide over time helps identify patients with persistent endogenous insulin secretion (brittle vs. stable diabetes) and informs insulin requirements. Glucagon-stimulated C-peptide >0.2 nmol/L correlates with improved HbA1c and fewer hypoglycemic events
  \item Evaluation of fasting hypoglycemia: a C-peptide level that is inappropriately high (>0.6 nmol/L) with simultaneous glucose <70 mg/dL suggests endogenous hyperinsulinemia (insulinoma, sulfonylurea or meglitinide use). Conversely, a low C-peptide (<0.2 nmol/L) with hypoglycemia and a high serum insulin suggests exogenous insulin administration (factitious hypoglycemia) or insulin antibody-mediated hypoglycemia
  \item Monitoring of beta-cell function after islet cell transplantation and total pancreatectomy with islet autotransplantation (TPIAT): C-peptide >0.3 nmol/L is the primary marker of graft function
  \item Assessment of beta-cell function in cystic fibrosis-related diabetes (CFRD), where annual screening with OGTT plus C-peptide is sometimes used
  \item Research and clinical trials evaluating beta-cell preservation therapies (immunomodulators, teplizumab) in recent-onset type 1 diabetes
\end{itemize}

\section*{Primary vs Secondary Test}
C-peptide is a primary test for differentiating diabetes type when the clinical picture is atypical (age >30, obesity, or negative autoantibodies) and for evaluating unexplained fasting hypoglycemia. It is a secondary/supplementary test in established diabetes, where its main role is assessing residual beta-cell function and guiding intensification of insulin therapy.

\section*{How This Test Is Performed}
A healthcare professional draws blood from a vein. The preferred sample is fasting (8--12 hours), collected in a serum separator or EDTA tube (lavender-top). The sample is centrifuged and separated promptly, then analyzed by immunoassay. Some laboratories offer the glucagon stimulation test: after a baseline fasting C-peptide sample, 1 mg of glucagon is administered IV, and a second C-peptide sample is drawn exactly 6 minutes later. For the 24-hour urine C-peptide test, all urine is collected over 24 hours in a container (refrigerated); a C-peptide-to-creatinine ratio is calculated.

\section*{What Preparation Is Needed}
Fasting (8--12 hours) is required for the fasting C-peptide test. Water is permitted. Avoid strenuous exercise the morning before the test (it can alter insulin secretion). Medications that affect insulin secretion (sulfonylureas, meglitinides, GLP-1 receptor agonists, corticosteroids, diazoxide) should be documented; the provider will decide whether to withhold them. Do not stop insulin or diabetes medications without instructions. For the glucagon-stimulated test, glucagon is contraindicated in patients with pheochromocytoma and in those with suspected insulinoma (may cause severe hypoglycemia due to prolonged insulin secretion).

\section*{Contraindications and Precautions}
\begin{itemize}
  \item C-peptide is not valid for diagnosing diabetes type if the patient is receiving exogenous insulin and has insulin antibodies or if the specimen is hemolyzed (hemolysis falsely lowers C-peptide)
  \item C-peptide must ALWAYS be interpreted with simultaneous plasma glucose: an absolute value alone is meaningless without knowing the glucose. A C-peptide of 1.0 nmol/L is appropriate at glucose 250 mg/dL, but inappropriately high at glucose 60 mg/dL
  \item Sulfonylurea- and meglitinide-induced hypoglycemia produce HIGH C-peptide and HIGH insulin levels (because these drugs stimulate endogenous secretion). Measurement of sulfonylurea levels in the urine or serum can distinguish them from insulinoma. Insulin autoantibodies (IAA) may rarely cause false elevation of insulin assay results, complicating interpretation
  \item C-peptide does not cross-react with exogenous insulin in insulin assays, but insulin assays can cross-react with proinsulin, and proinsulin is elevated in some conditions (PC1/3 deficiency, familial hyperproinsulinemia), causing false elevation of insulin without a corresponding rise in C-peptide. This matters in the evaluation of hypoglycemia: the insulin-to-C-peptide ratio is helpful (high ratio >1 suggests exogenous insulin, since insulin and C-peptide are cleared differently)
  \item In renal insufficiency (eGFR <30 mL/min/1.73 m$^2$), C-peptide is cleared by the kidneys and may be falsely elevated (accumulation). Interpretation requires caution
  \item Urine C-peptide may be falsely high in renal impairment and falsely low in dilute urine; creatinine correction (C-peptide-to-creatinine ratio) partially compensates for this
\end{itemize}
\section*{Risks}
Standard venipuncture risks: minor pain, bruising, or hematoma. The glucagon stimulation test carries risks of nausea, vomiting, and transient hypoglycemia. In patients with insulinoma, glucagon can precipitate severe hypoglycemia. In patients with pheochromocytoma, glucagon can precipitate a hypertensive crisis. Both groups are excluded from the test or tested only under close monitoring.

\section*{What Happens After the Test}
Results are available within 1--3 days. The provider interprets the C-peptide with the simultaneous fasting glucose. In a patient with an elevated fasting glucose, a low C-peptide (<0.2 nmol/L) indicates severe insulin deficiency and warrants insulin therapy regardless of diabetes type. A fasting C-peptide >0.3 nmol/L generally indicates adequate residual beta-cell function to consider non-insulin therapies or to guide the choice of insulin regimens (basal-only vs. basal-bolus). In a patient with fasting hypoglycemia, a high C-peptide with high insulin suggests endogenous hyperinsulinemia and prompts imaging (CT/MRI/endoscopic ultrasound) and sulfonylurea screening.

\section*{Understanding Results}

\subsection*{Below Normal}
\begin{itemize}
  \item \textbf{Very low or undetectable C-peptide (<0.2 nmol/L) with hyperglycemia:} Severe insulin deficiency. Consistent with type 1 diabetes (autoimmune destruction), LADA, or end-stage beta-cell failure in long-standing type 2 diabetes. Insulin therapy is usually required. Confirm with diabetes-related autoantibodies (GAD65, IA-2, ZnT8, IAA)
  \item \textbf{Low C-peptide with hypoglycemia:} Suggests exogenous insulin administration (factitious hypoglycemia), insulin antibody-mediated hypoglycemia, or, rarely, severe starvation with depleted hepatic glycogen. The serum insulin is elevated but C-peptide is suppressed because exogenous insulin suppresses endogenous beta-cell secretion
  \item \textbf{Low C-peptide in type 1 diabetes after a long duration:} Beta-cell function declines progressively. Loss of residual C-peptide (typically <0.1 nmol/L within 5--10 years) is associated with increased glycemic variability and higher risk of severe hypoglycemia
  \item \textbf{Low C-peptide with normal glucose:} May reflect physiologic suppression (prolonged fasting) or beta-cell dysfunction without overt hyperglycemia. Interpret with caution; a normal fasting glucose suggests preserved beta-cell function overall
\end{itemize}

\subsection*{Above Normal}
\begin{itemize}
  \item \textbf{Inappropriately high C-peptide (>0.6 nmol/L) with fasting hypoglycemia (<70 mg/dL):} Endogenous hyperinsulinemia. Causes: insulinoma (beta-cell adenoma; 90\% benign, 10\% malignant). Up to 10\% are multiple (MEN1). Sulfonylurea or meglitinide overdose (factitious). Persistent hyperinsulinemic hypoglycemia of infancy (PHHI; congenital hyperinsulinism). Autoimmune hypoglycemia (insulin antibodies bind and release insulin unpredictably). Exogenous insulin does NOT raise C-peptide
  \item \textbf{High C-peptide with hyperglycemia:} Insulin resistance with adequate beta-cell compensation (early type 2 diabetes), obesity, Cushing syndrome, acromegaly, pregnancy. In type 2 diabetes, C-peptide is normal or high in early disease and declines as beta-cell failure progresses
  \item \textbf{High C-peptide with elevated insulin and negative sulfonylurea screen:} Insulinoma should be suspected. Localize with pancreatic protocol CT/MRI, endoscopic ultrasound with fine-needle aspiration, or somatostatin receptor scintigraphy (68Ga-DOTATATE PET/CT)
  \item \textbf{High C-peptide in renal insufficiency:} Decreased renal clearance causes C-peptide accumulation. This is expected and NOT a marker of hyperinsulinemia
\end{itemize}

\section*{Normal Results}
\begin{itemize}
  \item \textbf{Fasting plasma C-peptide (normoglycemic, euglycemic adults):} 0.5--2.0 ng/mL (0.17--0.66 nmol/L); typical values 1.1--4.4 ng/mL (0.37--1.47 nmol/L) in some laboratories
  \item \textbf{Stimulated (glucagon or mixed-meal) C-peptide:} usually >2 ng/mL (0.66 nmol/L) and often 4--8 ng/mL (1.3--2.6 nmol/L) in healthy individuals
  \item \textbf{Urine C-peptide:} 10--120 $\mu$g/day in adults; urine C-peptide-to-creatinine ratio typically <4 nmol/mmol creatinine
  \item \textbf{Reference ranges vary by assay and laboratory; always use the laboratory-specific range on the report}
\end{itemize}

\section*{Follow-up Tests}
\begin{itemize}
  \item \textbf{Diabetes autoantibody panel:} GAD65, IA-2, ZnT8, IAA (islet cell antibodies) to confirm type 1 diabetes when C-peptide is low and the clinical picture is atypical (age >30, slow onset)
  \item \textbf{Simultaneous plasma glucose and serum insulin:} Required for interpretation of C-peptide. Calculate insulin-to-C-peptide ratio: a molar ratio >1 is highly suggestive of exogenous insulin (factitious), whereas a ratio <1 supports endogenous secretion
  \item \textbf{Sulfonylurea/meglitinide screen (urine or serum):} Distinguishes drug-induced hyperinsulinemia (sulfonylurea use) from insulinoma in patients with high C-peptide and hypoglycemia
  \item \textbf{Pancreatic imaging:} CT abdomen with pancreatic protocol, MRI/MRCP, endoscopic ultrasound with FNA, or 68Ga-DOTATATE PET/CT for insulinoma localization. In MEN1 screening, C-peptide is combined with prolactin, calcium, PTH, and imaging
  \item \textbf{72-hour fast test:} The gold standard for the diagnosis of insulinoma. Performed when C-peptide, insulin, proinsulin, and beta-hydroxybutyrate are measured every 4--6 hours until symptoms occur, glucose <45 mg/dL, or 72 hours elapse. Insulinoma is diagnosed when C-peptide is suppressed inappropriately and beta-hydroxybutyrate is low. This test requires inpatient monitoring
  \item \textbf{HbA1c:} For monitoring glycemic control in patients with diabetes after C-peptide-guided classification and treatment adjustment
\end{itemize}

\section*{References}
\bibliographystyle{unsrtnat}
\bibliography{references}

\clearpage
\section*{Regulatory Status and Authorization Date}
This chapter describes a test category, not a specific commercial product. FDA, CDSCO, EU IVDR, and MHRA approval or registration dates therefore cannot be assigned without the assay manufacturer, product name, instrument, and intended use. For laboratory-developed tests, an individual agency approval date may not exist. See the Preface for the interpretation of regulatory status.

\clearpage
